\begin{abstract}
LLM-based agents operate under natural-language policy specifications that inevitably contain ambiguity, yet the causal relationship between ambiguity type and safety violations remains unmeasured. We present a matched-pair study on $\tau^2$-bench with 300 clause pairs spanning a six-type ambiguity taxonomy, evaluated across five models (2,997 episodes). Three findings emerge. First, ambiguous clauses increase violation rates by 34.0 percentage points (OR $= 15.1$), with a perturbation control supporting ambiguity-specific causation. Second, incompleteness consistently ranks highest (46--62\% across seven models); the remaining five types cluster into a lower tier at the $\pm$15\,pp margin, though tier cohesion is margin-dependent. Third, authorization scope shows the widest but least reliably measured cross-model range (52\,pp; $\kappa = 0.385$, \S Limitations), suggesting that safety training may differentially mitigate specific ambiguity types. Agents show fewer than 11\% disambiguation awareness, adopting surface interpretations without flagging uncertainty. These results motivate type-aware policy auditing and externalized ambiguity detection for policy-governed agent deployments.
\end{abstract}
